\section{Evaluation}
\label{sec:evaluation}

\subsection{Experimental Setup}

\paragraph{Benchmarks.} We evaluate on two SV-COMP benchmark sets. The primary set is
the SV-COMP Loops ReachSafety suite (\texttt{loops\_reachsafety\_unreach}), comprising
764 tasks drawn from the official \texttt{Loops.set} with the
\texttt{unreach-call.prp} property specification. This set is chosen because loop
invariant discovery is precisely the bottleneck that VGuide targets, making it the most
informative evaluation domain. The secondary set is the NoOverflow scalar subset
(\texttt{no\_overflow\_scalar}), comprising 452 tasks, used in the competition-grade
combined evaluation to assess cross-domain transfer.

\paragraph{Time limit.} All experiments use a 300\,s per-task wall-clock time limit,
matching the official SV-COMP competition setting. Memory limit is 15\,GB per task.

\paragraph{Tool and configurations.} We use CPAChecker with VGuide extensions. Two
configurations are evaluated. \texttt{predicateAnalysis-vguide} is the base predicate
analysis configuration with VGuide enabled; its stock baseline is standard
\texttt{predicateAnalysis} without VGuide. \texttt{svcomp26-vguide} is the full
SV-COMP 2026 competition portfolio with VGuide integrated into the predicate-analysis
component; its stock baseline is \texttt{svcomp26} without VGuide. The portfolio
configuration runs multiple analysis strategies in parallel---symbolic execution,
k-induction, value analysis, and predicate analysis---with VGuide active only in the
predicate-analysis worker.

\paragraph{LLM.} All VGuide runs use \LLMName{} (non-thinking mode). The LLM is queried
once per task, on the first spurious CE, with the dual-prompt design described in
\Cref{sec:design}. API access uses the standard DeepSeek endpoint with a fixed
temperature of~0 for reproducibility.

\paragraph{Metrics.} We report (i)~tasks solved, counting a task as solved if
CPAChecker returns \textsc{True} or \textsc{False} (i.e., a verdict is produced, not
\textsc{Unknown} or an error); (ii)~PAR-2 score, defined as the average of wall-clock
time for solved tasks and $2 \times 300 = 600$\,s for unsolved tasks, lower being
better; and (iii)~wrong verdicts, the count of tasks for which the tool returned
\textsc{True} on an expected-\textsc{False} task or vice versa.

\subsection{Main Results: Loops ReachSafety}

\Cref{tab:loops} reports results on the 764-task Loops ReachSafety set for all
four configurations.

\begin{table}[t]
  \caption{Results on SV-COMP Loops ReachSafety (764 tasks, 300\,s limit).
           \emph{Solved}: tasks with a definite verdict. \emph{PAR-2}: lower is
           better. \emph{Wrong}: wrong verdicts (soundness violations).}
  \label{tab:loops}
  \centering
  \begin{tabular}{lccc}
    \toprule
    Configuration & Solved & PAR-2 (s) & Wrong \\
    \midrule
    predicateAnalysis (stock)   & \LoopsStock{}          & \LoopsPARStock{}          & 0 \\
    predicateAnalysis + VGuide  & \LoopsVGuide{}         & \LoopsPARVGuide{}         & 0 \\
    \addlinespace
    svcomp26 (stock)            & \PortfolioLoopsStock{} & \PortfolioPARStock{}      & 0 \\
    svcomp26 + VGuide           & \PortfolioLoopsVGuide{}& \PortfolioPARVGuide{}     & 0 \\
    \bottomrule
  \end{tabular}
\end{table}

VGuide gains \LoopsDelta{} tasks over the base predicate analysis
(\LoopsStock{}$\to$\LoopsVGuide{}), a 16.4\% relative improvement, and reduces PAR-2
by \LoopsPARDelta{}\,s (6.2\% reduction). The performance improvement is achieved with
zero soundness violations across all \LoopsVGuide{} solved tasks and all unsolved tasks.

On the full svcomp26 portfolio, the gain is \PortfolioLoopsDelta{} tasks
(\PortfolioLoopsStock{}$\to$\PortfolioLoopsVGuide{}), with PAR-2 reduced by
\PortfolioPARDelta{}\,s. The smaller absolute gain is expected and arises from
portfolio saturation: the svcomp26 baseline already deploys symbolic execution,
k-induction, and value analysis in parallel. Tasks that VGuide can uniquely solve on
the base predicate configuration are frequently solved by one of these other strategies
in the portfolio, leaving fewer tasks where VGuide's predicate hints provide the sole
path to a verdict. The \PortfolioLoopsDelta{} gain therefore represents cases where
VGuide's predicate injection enables the predicate-analysis worker to return a result
faster than any parallel worker could.

All \LoopsDelta{} incremental wins in the base configuration are \textsc{True} verdicts
(safety proofs); the \textsc{False} count is unchanged. This is consistent with the
design: VGuide injects predicates that support invariant proofs; it does not provide
counterexample witnesses for genuine bugs.

\subsection{Competition-Grade Combined Results}

\Cref{tab:combined} reports the competition-grade combined evaluation merging
Loops ReachSafety and NoOverflow scalar results under the svcomp26 portfolio
configurations.

\begin{table}[t]
  \caption{Competition-grade combined results (ReachSafety + NoOverflow scalar,
           300\,s limit, svcomp26 portfolio). \emph{$\Delta$}: gain over stock.}
  \label{tab:combined}
  \centering
  \begin{tabular}{lcccc}
    \toprule
    Category & Stock & VGuide & $\Delta$ & PAR-2 (stock / VGuide) \\
    \midrule
    ReachSafety  & \CompReachStock{} & \CompReachVGuide{} & \CompReachDelta{} & \CompPARReachStock{}\,s / \CompPARReachVGuide{}\,s \\
    NoOverflow   & \CompOvfStock{}   & \CompOvfVGuide{}   & \CompOvfDelta{}   & \CompPAROvfStock{}\,s  / \CompPAROvfVGuide{}\,s  \\
    \midrule
    Total        & ---               & ---                & \CompTotalDelta{} & --- \\
    \bottomrule
  \end{tabular}
\end{table}

VGuide improves on both categories: \CompReachDelta{} additional ReachSafety tasks
(\CompReachStock{}$\to$\CompReachVGuide{}) and \CompOvfDelta{} additional NoOverflow
tasks (\CompOvfStock{}$\to$\CompOvfVGuide{}). PAR-2 drops by \CompPARReachDelta{}\,s
on ReachSafety and by \CompPAROvfDelta{}\,s on NoOverflow. No wrong verdicts are
produced on either category.

The transfer to NoOverflow is noteworthy: VGuide's context pack and validation cascade
were designed with ReachSafety in mind, yet the predicate injection strategy extends to
overflow-checking tasks without modification. Overflow tasks share the loop counting
patterns of the ReachSafety set---arithmetic on bounded integer variables---and the
LLM's predicate suggestions are similarly relevant.

\subsection{Ablation: Is CE Context Necessary?}

To isolate the contribution of CE-context from static program understanding, we
compare three configurations on the base predicate analysis setting.

\begin{table}[t]
  \caption{Ablation study on SV-COMP Loops ReachSafety (764 tasks, 300\,s).
           \emph{source\_prior}: LLM fires at analysis start with source only, no
           CE trace. \emph{first\_spurious}: LLM fires on first spurious CE (full
           VGuide).}
  \label{tab:ablation}
  \centering
  \begin{tabular}{lccc}
    \toprule
    Configuration & Solved & PAR-2 (s) & Wrong \\
    \midrule
    predicateAnalysis (stock) & \LoopsStock{}           & \LoopsPARStock{}  & 0 \\
    source\_prior             & \AblationSourcePrior{}  & 427.1             & 0 \\
    first\_spurious (VGuide)  & \AblationFirstSpurious{}& \LoopsPARVGuide{} & 0 \\
    \bottomrule
  \end{tabular}
\end{table}

The \emph{source\_prior} variant fires the LLM at the beginning of the analysis, before
any CE has been found, with the stripped C source but no CE trace or block formula. The
LLM proposes predicates based purely on static program structure; these are validated
through the same L1/L2/L3 cascade and injected into the initial precision.

The result is unambiguous: source\_prior solves \AblationSourcePrior{} tasks, identical
to the stock result, with PAR-2 marginally worse (427.1\,s vs.\ \LoopsPARStock{}\,s).
The initial predicate injection provides no benefit whatsoever. In contrast,
first\_spurious (full VGuide) solves \AblationFirstSpurious{} tasks (\LoopsDelta{}
over stock).

This ablation establishes a clear causal claim: the \LLMName{} predicate oracle is
entirely without value when given only source code. Its contribution is activated
exclusively by the CE trace. The LLM's capability is not \emph{static pattern recognition
on source structure} but \emph{relational reasoning over concrete execution paths}. This
finding validates the central design decision of VGuide: CE-guided refinement is the
correct and necessary integration point for LLM predicates in a CEGAR-based verifier.
